\appendix
\chapter{Implementation Parameters}
\label{app:impl_params}

This appendix documents the detailed implementation parameters of the tactic implementation subsystem, which were condensed from the main text for readability.

\section{BM25 Retrieval}
\label{sec:app_bm25}

File retrieval uses the BM25 ranking algorithm with $k_1 = 1.5$ and $b = 0.75$, following standard information retrieval defaults~\cite{robertson2009probabilistic}. The query is constructed from the tactic name, category, and description. Each file's document is formed by concatenating its path, class names, function names, and import list. The top $k = 5$ files with positive BM25 scores are passed to the Planner Agent.

\section{Iteration Control}

The agentic loop terminates when any of the following conditions is met:

\begin{itemize}
    \item Maximum iterations reached (default: 5);
    \item Planner Agent issues a \texttt{STOP} signal;
    \item Planner exceeds 3 consecutive parse or validation failures.
\end{itemize}

\section{Patch Generation Constraints}

The Patch Agent returns the full replacement file content (not a diff). Generated files are rejected if they exceed 400 lines or target \texttt{\_\_init\_\_.py}. The agent is instructed to make one focused change per step and to preserve existing external interfaces.

\section{Self-Healing}

When a regression is detected, the affected file is restored from the pre-patch backup and the Patch Agent is re-invoked with test failure output as feedback. Up to 2 repair attempts are permitted per iteration. If all attempts fail, the iteration is abandoned and the pipeline continues to the next iteration.

\section{Backup Directory Structure}

File-level backups are organized per iteration and attempt:

\begin{verbatim}
backups/<repository>/step_<N>/attempt_<M>/<file_path>
\end{verbatim}

This structure enables precise rollback of individual files without reverting unrelated changes from other iterations.

\section{Artifact Storage}

For each repository, the following artifacts are preserved:

\begin{itemize}
    \item \textbf{Structural index:} JSON file with per-file path, imports, class names, function names, and size.
    \item \textbf{Test baseline:} pre-modification pytest output with failure list.
    \item \textbf{Step artifacts:} one JSON patch per completed iteration.
    \item \textbf{Test results:} pytest output, failure names, regression flag, and success status per iteration.
    \item \textbf{Artifact extraction:} file tree, import graph (text and JSON), code signatures.
\end{itemize}

\section{Pipeline Execution Workflow}

Per repository, the system executes sequentially:

\begin{enumerate}
    \item Shallow clone repository;
    \item Extract file tree, import graph, code signatures, structural metrics;
    \item Compute baseline MI, fan-out, package count, docstring coverage;
    \item Run LLM-based architecture detection;
    \item Run LLM-based tactic selection;
    \item Apply LLM-generated modifications via agentic loop;
    \item Validate syntactic correctness of modified files;
    \item Recompute all metrics on modified repository;
    \item Store results and logs;
    \item Delete local repository clone.
\end{enumerate}
